% !PRINTED 352
\weberpage{352}

%% [running head, not body text -- commented out 2026-08-07] Page furniture (running
%% head, folio, signature) is excluded from this edition by standing policy.
%% \noindent 352\hfill Zwölfter Abschnitt.\hfill \S\ 96.

If $\Delta$ is the fundamental number of the field $\Omega$, we put, according as $\Delta$ is even or odd:
\weq{1}{
\begin{aligned}
\theta_0 &= \tfrac{1}{2}\sqrt{\Delta}, & &\text{or}\\
\theta_0 &= \frac{-1+\sqrt{\Delta}}{2}\,.
\end{aligned}
}

Then by \S\ 90 every integer of the field $\Omega$ is representable in the form
\weq{2}{
\omega = x + y\theta_0
}
where $x$ and $y$ are integers, and this number is then and only then congruent to a rational number $a$ with respect to the modulus $Q$, when $y$ is divisible by $Q$, because then
\[
\frac{x - a + y\theta_0}{Q}
\]
must be an integer. Accordingly all numbers of the order $[Q]$ are contained in the form:
\[
\omega = x + yQ\theta_0,
\]
or, in the two cases distinguished in (1), if
\weq{3}{
D = Q^2\Delta
}
is put:
\weq{4}{
\begin{aligned}
&x - y\,\frac{\sqrt{D}}{2}, & &\Delta \text{ even},\\
&x - \frac{Q}{2}y + y\,\frac{\sqrt{D}}{2}, & &\Delta \text{ odd}, D \text{ even},\\
&x - \frac{Q-1}{2}y + y\,\frac{-1+\sqrt{D}}{2}, & &D \text{ odd}.
\end{aligned}
}

If we therefore put
\weq{5}{
\begin{aligned}
\theta &= \frac{\sqrt{D}}{2}, & D &\equiv 0 \pmod{4}\\
\theta &= \frac{-1+\sqrt{D}}{2}, & D &\equiv 1 \pmod{4}
\end{aligned}
}
and in the last two formulae (4) replace the arbitrary integers
\[
x - \frac{Q}{2}y, \quad x - \frac{Q-1}{2}y
\]
by $x$, we obtain the result:
%CONT>>
